\section{Comparing the marked local integrals}
\label{sec:marked-kernel-convergence}
We now supply the local analytic comparison used in
Section~\ref{sec:local-source-stability}. Fix an adjacent pair of
scales and one period. Compare two sequences of paid transitions
on the same support dictionary. Assume their fine and coarse
canonical path coefficients, retained Gaussian rows and couplings
have common limits. The path sums have the strict-width bounds
proved above. For the two backgrounds assume that each ordinary
and forced activity difference tends uniformly to zero and that
their activity majorants have uniformly small tails. In a marked
comparison the inserted function and its address are the same
on both sides and its paid size is bounded. These assumptions
will later be supplied by the background trajectory construction.

The claim is convergence of the difference of each local marked
integral and of every paid derivative on its healthy support.
It includes uniformity over moving coarse configurations and
normalized local charts. The principal issue is that the Gaussian
measures in the two integrals vary. We first put them on a common
physical variable; this also keeps possibly discontinuous inserted
functions at exactly the same argument.

\subsection{One common physical integration variable}
Fix a forest, its finite coordinate choice and the support pattern
of a local term. Its integral is over a compact gauge slice $Y$
with probability Haar measure and a real Gaussian coordinate
$z\in\mathbb R^D$. The coordinate change in the local integral
has the form
\[
                     w=\mu_N(Y)+g_N z,\qquad
                     z=r_N(w-\mu_N(Y)),\quad r_N=g_N^{-1}.
\]
Coordinate permutations belonging to the fixed support pattern
are included in this notation; they have unit absolute Jacobian.
The paid coupling window gives constants $0<u\le v<\infty$
such that $u\le r_N\le v$, and the local mean estimate gives
$\|\mu_N(Y)\|\le B$ independently of $N,Y$. If $\rho_N$ is
the centered Gaussian density, the transformed density is
\[
             d_N(Y,w)=r_N^D\rho_N(r_N(w-\mu_N(Y))).
\]
The actual covariance bounds give fixed $C,\tau>0$ with
$\rho_N(z)\le C e^{-\tau\|z\|^2}$. Since
\[
 \|w\|^2\le2\|w-\mu_N(Y)\|^2+2B^2,
\]
we have
\[
 d_N(Y,w)\le v^DCe^{\tau u^2B^2}
                              e^{-\tau u^2\|w\|^2/2}.
\tag{$\ddagger$}\label{eq:marked-common-gaussian-envelope}
\]
Also
\[
 1+\|r_N(w-\mu_N(Y))\|
                  \le(1+v+vB)(1+\|w\|).
\]
Thus a local integrand with bound $C_k(1+\|z\|)^k$ becomes
bounded, after multiplication by $d_N$, by a constant times
\[
                  e^{-\tau u^2\|w\|^2/2}(1+\|w\|)^k.
\]
This is integrable on $\mathbb R^D$. The constants may depend
on the fixed scale, period, support and derivative order, but
not on the terminal cutoff, the slice point, or the moving
coarse configuration.

The local differentiation theorem of
Section~\ref{sec:gaussian-differentiation} writes each derivative
as an integral of the differentiated compensated integrand
against the Gaussian measure frozen at the base configuration.
The bounds in Sections~\ref{sec:packet-estimates}
and~\ref{sec:marked-activity-jets} give the polynomial bound
just used. Therefore the change of variable is applied to that
already justified derivative integral; it does not differentiate
a moving measure without compensation.

\subsection{Why the transformed integrand differences vanish}
First hold the finite support pattern and coordinate choice fixed.
At a fixed physical point $(Y,w)$ the fine configuration read by
an inserted activity is independent of the Gaussian covariance
and of the coupling: the identity
$\mu_N+g_Nr_N(w-\mu_N)=w$ cancels their dependence exactly.
Hence a common bounded marked function is evaluated at the same
fine configuration on both sides. No continuity of that function
is needed. An unmarked background activity has the same property;
its two values differ by the assumed uniform activity difference.

The other factors in the compensated integrand are smooth finite
expressions in the local path and retained-kernel data. Their
continuity can be checked directly. A fixed length cutoff leaves
only finitely many canonical paths: their displacement vectors
lie in a finite lattice box. Each path holonomy is a finite
product of link matrices; all its fixed-order derivatives are
finite sums of such products. Convergence of its coefficient
therefore gives convergence of these derivatives along any
sequence of bounded local chart jets. Infinite canonical path
sums have the same conclusion by the strict-width summability
proved in Section~\ref{sec:fixed-rung-convergence}.

The local Gram matrices are finite sums and products of those
rows. The common positive lower bound passes to their limits.
Their inverses converge by
\[
              A^{-1}-B^{-1}=A^{-1}(B-A)B^{-1}.
\]
Repeated differentiation of $A A^{-1}=I$ proves convergence
of all required inverse jets by induction. Determinants are
polynomials in the entries, and the determinant stays strictly
positive. Their logarithms and inverse square roots, and their
fixed-order derivatives, consequently converge. This covers
the selected covariance, means, Gaussian densities and density
ratios, including the projected covariance in the compensation.
The denominator density at a fixed $w$ has a strictly positive
limit; its reciprocal therefore causes no loss of pointwise
convergence. Uniform domination is supplied separately by
\eqref{eq:marked-common-gaussian-envelope} and the compensated
derivative bounds.

Before differentiating these factors, retain the full protected
support of each activity and of the marked insertion. The base coarse
configuration is healthy on the enlarged union of these supports,
and the normalized chart is based there. Equation~\eqref{eq:compensator-held-fine}
holds on every active protected bond wherever the broader cutoff
is nonzero near the base point; inactive bonds stay fixed by definition.
Indeed, a nonzero cutoff germ gives the base fluctuation bound $4at$;
continuity gives the strict $8at$ bound on a neighborhood. The
compensator identity then applies on all the finitely many protected
bonds on one common neighborhood. Where a cutoff has zero germ,
the complete factor has zero germ. Consequently an activity or mark
$h_N$ occurs, as a germ in the chart parameter $y$, only as
\[
                         h_N(U_N(Y,w))\,\zeta_N(y;Y,w),
\]
with $U_N(Y,w)$ independent of $y$ and $\zeta_N$ smooth. This is
exactly the held-factor identity~\eqref{eq:marked-held-factor},
applied to each relevant read support. All derivatives fall on
$\zeta_N$ and the remaining smooth factors. Uniform convergence
of activity values is therefore sufficient; convergence of their
uncompensated derivatives is not being inferred.

For completeness, the elementary factors in a local product fall
into four mathematical cases. A cell-action factor is the
exponential of a finite local action and fixed cutoff functions;
its action jets converge by the preceding path argument. A
forced-activity factor is its value at the common fine configuration
times a smooth local cutoff. An ordinary logarithmic factor is
the exponential minus one of an ordinary activity value together
with the local quadratic compensation; the first value is
compared uniformly and the second by the matrix argument. A
warning factor is a fixed local mask or charge times a smooth
cutoff, possibly minus its constant term. The flag and compatibility
decisions have the same values on both sides, because their
geometric support data, threshold and physical configuration
are identical. Product and chain rules now prove that the differences
of the fixed-order jets of every finite product tend to zero.
With one marked factor, replace the corresponding ordinary
factor by the common inserted function times its cutoff; the
same proof applies.

Moving configurations do not require convergence along the
whole sequence. For any subsequence, take a further subsequence
on which their finitely many link matrices and the bounded
chart jets converge. An order-$r$ derivative evaluation can be
pulled back to $r$ real parameters, so these finitely many jets
lie in one fixed finite-dimensional bounded set. Compactness of
the gauge group and of the closure of that set gives the subsequence. The support and
coordinate choices range over finite sets, so take them constant
on a further subsequence. The calculations above make the difference of the two
transformed integrands tend to zero there. A common bounded
mark need not itself converge: it multiplies a smooth-factor
difference tending to zero, and its uniform bound suffices. If a difference
failed to tend to zero in the original sequence, a subsequence
would stay a positive distance from zero and could not have
such a further subsequence. Thus the differences tend to zero
without any preliminary convergence assumption on the moving
configurations. Shared discontinuous indicators may also vary: for their common
value $\mathbf1_{E_N}$ use
$|\mathbf1_{E_N}(A_N-B_N)|\le|A_N-B_N|$. Only the finite
structural support and coordinate choices are frozen; no
simultaneous stabilization of indicators for every $(Y,w)$ is needed.

\subsection{Integration and local derivative norms}
Let $F_N(Y,z)$ and $G_N(Y,z)$ be the two compensated derivative
integrands. For fixed $Y,w$, their density-weighted transforms
have difference tending to zero. Each is dominated by an
integrable function of $w$ independent of $Y,N$. Dominated
convergence first in $w$ proves convergence of the difference
of their inner integrals for almost every $Y$. The norms of
these differences are bounded by the integral of the sum of
the two majorants. A second application of dominated convergence
against the probability Haar measure in $Y$ gives convergence
of the complete local integrals. This argument also justifies
the two uses of Fubini, since absolute integrability was established
before passing to the limit.

Apply this to every moving base configuration and normalized
chart of each order up to the paid order. The selection argument
of Section~\ref{sec:local-source-stability} turns it into
uniform convergence of those derivative evaluations. On the
fixed finite torus, the conversion from normalized chart
derivatives to the local derivative price has a finite constant:
all site separations are bounded, and only finitely many derivative
orders occur. Therefore each individual ordinary output entry
has difference tending to zero in its required derivative price.
For the forced entry use order zero and arbitrary moving
configurations; its uniform value difference tends to zero.

Finite averaging over forests and coordinate choices preserves
both statements. The ordinary marked logarithm is a sum of
finite products with one marked factor, and the forced marked
polynomial is formed by finite-subset products, addition of the
ordinary logarithm, and differentiation in its one source
parameter. Each operation preserves the entrywise convergence
just proved. The marked cluster and product estimates in
Section~\ref{sec:local-source-stability} give their uniform
tails. Local subtraction and the background hook then give
the convergence of each fixed input column in the actual
common dictionary. Summation and finite motion averaging give
the adjacent-scale stability required in
Section~\ref{sec:fixed-rung-convergence}.

This comparison uses both backgrounds' common-limit hypotheses;
it does not prove convergence of the background activity trajectory.
It proves the needed continuous dependence of the selected
ordinary and marked local transformations, with their masks,
Gaussian measures, derivatives and countable output lists retained.
